\documentclass[11pt]{article}
\usepackage[margin=1in]{geometry}
\usepackage{amsmath,amssymb}
\usepackage{enumitem}
\usepackage[T1]{fontenc}
\usepackage[utf8]{inputenc}
\usepackage{lmodern}
\setlength{\parindent}{0pt}
\setlength{\parskip}{0.5em}
\begin{document}

\section*{IB Physics HL: D2 + D3 Problem Collection}
Paraphrased from 2025-syllabus past papers; mixed questions are included when they materially assess electric/magnetic fields or motion in those fields.
\begin{itemize}[leftmargin=1.25em]
\item
\textbf{Source.} May 2025 HL Paper 1A TZ2, Question 25\\
\textit{Original format: multiple choice.}\\
In a uniform magnetic field directed into the page, three incoming particles follow three different paths labelled $Q$, $R$, and $S$. Determine which path corresponds to each particle type (proton, neutron, electron, or alpha particle), using the curvature and direction of motion.
\item
\textbf{Source.} May 2025 HL Paper 1A TZ2, Question 26\\
\textit{Original format: multiple choice.}\\
A cell, two parallel rails, and a movable conducting rod of length $L$ form a circuit. A uniform magnetic field points out of the page, and the rod moves to the right with speed $v$. Determine the direction of the current in the circuit and the initial magnetic force on the rod.
\item
\textbf{Source.} May 2025 HL Paper 1A TZ2, Question 27\\
\textit{Original format: multiple choice.}\\
Two long parallel wires $X$ and $Y$ initially carry equal currents and repel with force per unit length $F$. Then both currents are reversed, the current in $X$ is doubled, the current in $Y$ is unchanged, and the separation is doubled. Determine the new force per unit length and whether it is attractive or repulsive.
\item
\textbf{Source.} May 2025 HL Paper 1A TZ3, Question 29\\
\textit{Original format: multiple choice.}\\
In a Millikan oil-drop setup, one drop is stationary between parallel plates. A second drop of the same density is placed between the same plates, but it has twice the charge and twice the radius of the first drop. Determine the initial motion of the second drop.
\item
\textbf{Source.} May 2025 HL Paper 1A TZ3, Question 30\\
\textit{Original format: multiple choice.}\\
Two isolated point charges, $+Q$ and $+2Q$, are separated by a distance $3d$. Point $P$ lies between them, at distance $d$ from the first charge and $2d$ from the second. Determine the net electric field strength at $P$.
\item
\textbf{Source.} May 2025 HL Paper 1A TZ3, Question 31\\
\textit{Original format: multiple choice.}\\
Two positive charges $Q_1$ and $Q_2$ are placed symmetrically around a third charge $Q_3$. The force on $Q_3$ due to $Q_1$ has magnitude $F$, while the force due to $Q_2$ has magnitude $2F$. Determine the resultant force on $Q_3$.
\item
\textbf{Source.} November 2025 HL Paper 1A TZ3, Question 30\\
\textit{Original format: multiple choice.}\\
A positively charged rod is brought near a grounded metal plate. The grounding wire is removed first, and then the rod is removed. Determine the net charge on the plate immediately before the grounding wire is removed and after the rod is removed.
\item
\textbf{Source.} November 2025 HL Paper 1A TZ3, Question 31\\
\textit{Original format: multiple choice.}\\
An electron enters midway between two charged parallel plates with speed $v$. The potential difference between the plates is $V$. Determine the electron's speed when it strikes one of the plates.
\item
\textbf{Source.} November 2025 HL Paper 1A TZ3, Question 32\\
\textit{Original format: multiple choice.}\\
A positively charged particle moves toward a straight current-carrying wire. Determine the direction of the magnetic force acting on the particle.
\item
\textbf{Source.} May 2025 HL Paper 2 TZ2, Question 3\\
Two oppositely charged parallel plates are separated by $8.0\,\text{cm}$ and have a potential difference of $120\,\text{V}$. An alpha particle is released from rest at the positive plate. Ignore gravity.

\begin{enumerate}[label=(\alph*), leftmargin=1.5em]
\item Calculate the electric field strength between the plates. [1]
\item For the alpha particle:
\begin{enumerate}[label=(\roman*), leftmargin=1.5em]
\item show that its acceleration is about $7\times 10^{10}\,\text{m s}^{-2}$; [2]
\item calculate the time it takes to reach the negative plate; [2]
\item state its kinetic energy on arrival, in eV. [1]
\end{enumerate}
\item A magnetic field into the page is added, and an alpha particle now enters with horizontal speed $5.0\times 10^5\,\text{m s}^{-1}$ and is undeflected. Find the magnetic field magnitude. [2]
\item Replace the alpha particle in part (c) with an electron entering with the same speed. Sketch the electron\'s path. [1]
\end{enumerate}
\item
\textbf{Source.} May 2025 HL Paper 2 TZ3, Question 8\\
\textit{Note.} Mixed D2/D3/D4 question; included here because parts (a), (b), and (e) are directly on magnetic forces and magnetic fields.\\
A rod $R$ is perpendicular to a uniform magnetic field of strength $0.50\,\text{T}$ directed into the page. The rod carries current $I_R=2.0\,\text{A}$.

\begin{enumerate}[label=(\alph*), leftmargin=1.5em]
\item Calculate the force per unit length on $R$ and identify its direction. [2+1]
\item The external field is removed and a second long wire $Z$, carrying current $I_Z$, is placed parallel to $R$.
\begin{enumerate}[label=(\roman*), leftmargin=1.5em]
\item Show that the magnetic field due to $R$ at distance $r$ is $B_R=\dfrac{\mu_0 I_R}{2\pi r}$. [2]
\item Derive the SI base units of $\mu_0$. [1]
\end{enumerate}
\item Later in the same question, an electron moves parallel to a rod $T$ at speed $9.0\times 10^5\,\text{m s}^{-1}$. When a switch is closed, the magnetic field at the electron is $4.0\times 10^{-6}\,\text{T}$.
\begin{enumerate}[label=(\roman*), leftmargin=1.5em]
\item Calculate the initial radius of curvature of the electron\'s path. [3]
\item Explain how the radius subsequently changes. [2]
\end{enumerate}
\end{enumerate}
\item
\textbf{Source.} November 2025 HL Paper 2, Question 4\\
\textit{Note.} Mixed D3/D4 question; included here because it combines magnetic force on charges/currents with motional emf and Lenz-law reasoning.\\
A metal rod is pulled to the right at constant speed along two horizontal conducting rails $X$ and $Y$. The rails are connected by a straight conducting segment of length $L=0.60\,\text{m}$. The whole system is in a uniform magnetic field of magnitude $B=1.2\,\text{T}$ directed into the page.

\begin{enumerate}[label=(\alph*), leftmargin=1.5em]
\item The rod experiences a magnetic force of $0.084\,\text{N}$.
\begin{enumerate}[label=(\roman*), leftmargin=1.5em]
\item Show that the current in the rod is about $0.1\,\text{A}$. [1]
\item State the direction of the magnetic force on an electron in the rod. [1]
\item Explain why the magnetic field magnitude on the right side of the rod differs from that on the left side. [3]
\item Outline how Lenz\'s law applies in this situation. [2]
\end{enumerate}
\item Using the current from part (a)(i), determine the force per unit length between rails $X$ and $Y$, and state the SI base units of your answer. [3]
\end{enumerate}
\end{itemize}

\end{document}
